% Options for packages loaded elsewhere
\PassOptionsToPackage{unicode}{hyperref}
\PassOptionsToPackage{hyphens}{url}
\documentclass[
]{article}
\usepackage{xcolor}
\usepackage{amsmath,amssymb}
\setcounter{secnumdepth}{-\maxdimen} % remove section numbering
\usepackage{iftex}
\ifPDFTeX
  \usepackage[T1]{fontenc}
  \usepackage[utf8]{inputenc}
  \usepackage{textcomp} % provide euro and other symbols
\else % if luatex or xetex
  \usepackage{unicode-math} % this also loads fontspec
  \defaultfontfeatures{Scale=MatchLowercase}
  \defaultfontfeatures[\rmfamily]{Ligatures=TeX,Scale=1}
\fi
\usepackage{lmodern}
\ifPDFTeX\else
  % xetex/luatex font selection
\fi
% Use upquote if available, for straight quotes in verbatim environments
\IfFileExists{upquote.sty}{\usepackage{upquote}}{}
\IfFileExists{microtype.sty}{% use microtype if available
  \usepackage[]{microtype}
  \UseMicrotypeSet[protrusion]{basicmath} % disable protrusion for tt fonts
}{}
\makeatletter
\@ifundefined{KOMAClassName}{% if non-KOMA class
  \IfFileExists{parskip.sty}{%
    \usepackage{parskip}
  }{% else
    \setlength{\parindent}{0pt}
    \setlength{\parskip}{6pt plus 2pt minus 1pt}}
}{% if KOMA class
  \KOMAoptions{parskip=half}}
\makeatother
\setlength{\emergencystretch}{3em} % prevent overfull lines
\providecommand{\tightlist}{%
  \setlength{\itemsep}{0pt}\setlength{\parskip}{0pt}}
\usepackage{bookmark}
\IfFileExists{xurl.sty}{\usepackage{xurl}}{} % add URL line breaks if available
\urlstyle{same}
\hypersetup{
  pdftitle={The Trinity as a Necessary Structure of a Monistic Modal Ontology - Minimalist, Limit-Based Version},
  hidelinks,
  pdfcreator={LaTeX via pandoc}}

\title{The Trinity as a Necessary Structure of a Monistic Modal Ontology \\ \large Minimalist, Limit-Based Version}
\author{Paul Koop}
\date{}

\begin{document}
\maketitle

\begin{center}
\emph{Past and future are horizons of knowledge. \\ The present is the place of reality.}
\end{center}

\newpage
{
\setcounter{tocdepth}{3}
\tableofcontents
}
\newpage

\section{Preface}\label{preface}

This treatise is an attempt to demonstrate an ancient metaphysical intuition -- the threefold unity of origin, totality, and self-knowledge -- not as a matter of faith, but as a \textbf{logical necessity} of a consistent ontology.

It differs from earlier versions in three decisive respects:

\begin{enumerate}
\item \textbf{Minimal axiomatics:} Instead of six or nine axioms, the treatise will proceed with only two fundamental principles.
\item \textbf{Limit structure:} The concepts U (Origin), T (Totality), and S (Self-Knowledge) are not understood as substances, but as \textbf{limits} of an open interval.
\item \textbf{Transcendental deduction:} Axioms such as "consciousness is possible" or "modal realism" are not presupposed, but derived from monism and the existence of a world.
\end{enumerate}

The central thesis is:

> Under the axioms of monism and the existence of a possible world, a threefold limit structure of Origin, Totality, and Self-Knowledge necessarily follows.

The treatise provides \textbf{no proof of God} in the classical sense. It demonstrates no personal or substantial Trinity. It merely shows: \textbf{If you accept monism and the existence of a world, then a threefold limit structure necessarily follows.}

\section{Introduction: Aim and Methodological Self-Restriction}\label{introduction-aim-and-methodological-self-restriction}

Classical metaphysics has always attempted to derive the fundamental structure of reality from a few basic principles. A particular challenge arises from three seemingly distinct aspects:

\begin{enumerate}
\item \textbf{Why is there something rather than nothing?}
\item \textbf{How does the totality of all possibilities relate to actuality?}
\item \textbf{How can a state arise within reality that knows reality itself?}
\end{enumerate}

The ontology examined here proposes to answer these three questions not through three separate metaphysical substances, but through three necessary perspectives on the same reality:

\begin{itemize}
\item \textbf{Origin (U):} the necessary condition for the existence of possibilities at all – understood as the lower limit of Totality;
\item \textbf{Totality (T):} the entirety of all realized possibilities – understood as an open interval;
\item \textbf{Self-Knowledge (S):} the reflexive completion of reality in a state of complete knowledge of itself – understood as the upper limit of Totality.
\end{itemize}

This structure is called the \textbf{Trinity}:

\[
Tr := U \land T \land S
\]

The term "Trinity" here does not denote a personal or substantial threefoldness, but a functional unity of three necessary aspects of a single reality.

\section{Formal Language and Logical Framework}\label{formal-language-and-logical-framework}

\subsection{Modal Logic S5}\label{modal-logic-s5}

We work in first-order modal logic with identity in system S5:

\begin{itemize}
\item \textbf{Necessity:} \(\Box p\)
\item \textbf{Possibility:} \(\Diamond p := \neg\Box\neg p\)
\end{itemize}

\textbf{Axioms of S5:}
\begin{itemize}
\item (K) \(\Box(p \rightarrow q) \rightarrow (\Box p \rightarrow \Box q)\)
\item (T) \(\Box p \rightarrow p\)
\item (4) \(\Box p \rightarrow \Box\Box p\)
\item (5) \(\Diamond p \rightarrow \Box\Diamond p\)
\end{itemize}

\textbf{Inference Rules:}
\begin{itemize}
\item \textbf{Modus Ponens:} From \(p\) and \(p \rightarrow q\), infer \(q\).
\item \textbf{Necessitation:} From \(p\) (derivable), infer \(\Box p\).
\end{itemize}

\subsection{Predicate Logic}\label{predicate-logic}

We use classical first-order predicate logic with identity (\(=\)) and the usual quantifiers (\(\forall, \exists\)).

\section{The Two Axioms}\label{the-two-axioms}

\subsection{A1 – Monism}\label{a1-monism}

There are no fundamentally separated realms of reality.

\[
\boxed{A1 := \Box\neg\exists x\exists y\, FundamentallySeparated(x,y)}
\]

\textbf{Explanation:} If two realms were fundamentally separated, there could be no causal or ontological interaction between them. Then they could not be part of a unified reality, which contradicts monism. Such a dualism would be incompatible with a complete modal realism.

\subsection{A4 – Existence of a Possible Reality}\label{a4-existence-of-a-possible-reality}

There is at least one possible world.

\[
\boxed{A4 := \Diamond\exists w\, World(w)}
\]

\textbf{Explanation:} This axiom ensures that the modal universe is not empty. It is the weakest possible existence axiom: it does not say that a world \textbf{actually} exists, but only that it is \textbf{possible}.

\section{Definition of the Limit Structure}\label{definition-of-the-limit-structure}

\subsection{Totality T as an Open Interval}\label{totality-t-as-an-open-interval}

\textbf{Totality T} is the entirety of all realized possibilities. We understand T as an \textbf{open interval}:

\[
\boxed{T := (U, S)}
\]

This means: T is the set of all states that lie \textbf{between} Origin U and Self-Knowledge S. The limits U and S do \textbf{not} belong to T – they are \textbf{limit points}.

\textbf{What does "open interval" mean mathematically?}
\begin{itemize}
\item An open interval \((a, b)\) contains all numbers between a and b, but \textbf{not} a and b themselves.
\item It has \textbf{no smallest} and \textbf{no largest} element.
\item It has \textbf{limits} – a is the lower limit (infimum), b the upper limit (supremum).
\end{itemize}

\textbf{Transferred to ontology:}
\begin{itemize}
\item Totality T is the set of all \textbf{actually realized} states.
\item These states are \textbf{ordered} – from "minimal structure" to "maximal structure".
\item \textbf{Origin U} is the lower limit – that which comes closest to nothing, but is itself not nothing.
\item \textbf{Self-Knowledge S} is the upper limit – the complete transparency of Totality, which itself does not belong to Totality.
\end{itemize}

\subsection{Origin U as the Lower Limit}\label{origin-u-as-the-lower-limit}

\[
\boxed{U := \lim_{x \to \inf} T}
\]

This means: U is the \textbf{limit} of Totality T when approaching the "beginning." U is that which comes closest to nothing – but \textbf{not nothing} (for nothing would be a fundamental separation, which A1 prohibits).

\subsection{Self-Knowledge S as the Upper Limit}\label{self-knowledge-s-as-the-upper-limit}

\[
\boxed{S := \lim_{x \to \sup} T}
\]

This means: S is the \textbf{limit} of Totality T when approaching the "end." S is the complete self-transparency of Totality – but \textbf{not itself a part} of T (for otherwise it would not be complete).

\subsection{Well-Foundedness as a Consequence of Openness}\label{well-foundedness-as-a-consequence-of-openness}

From the openness of the interval follows \textbf{directly} well-foundedness:

\[
\boxed{\neg\exists infinite\, chain(g_1, g_2, ...)}
\]

\textbf{Justification:} In an open interval, there is no smallest element. Every element has a "before" – but there is \textbf{no final ground} that lies within the interval. The final ground is the \textbf{limit U}, which does \textbf{not} belong to the interval.

\subsection{The Trinity}\label{the-trinity}

The \textbf{Trinity} is the unity of the three limits:

\[
\boxed{Tr := U \land T \land S}
\]

Or in the language of limits:

\[
\boxed{Tr := \lim_{x \to \inf} T \;\land\; T \;\land\; \lim_{x \to \sup} T}
\]

\section{Derivation of Further Principles from A1 and A4}\label{derivation-of-further-principles-from-a1-and-a4}

\subsection{A5 – Consciousness as a Real Possibility (Theorem)}\label{a5-consciousness-as-a-real-possibility-theorem}

\textbf{Theorem:} From A1 and A4 follows \(\Diamond C\) (consciousness is possible).

\textbf{Proof:}

1. Suppose there is a world \(w\): \(\exists w\, World(w)\).

2. Suppose consciousness is impossible: \(\neg\Diamond C\). This means: \(\Box\neg C\) – it is necessary that there is no consciousness.

3. If there is no consciousness, then there is also no \textbf{experience} of world. The world would be \textbf{unknowable}.

4. An unknowable world would be \textbf{fundamentally separated} from any possible conscious perspective.

5. \textbf{But A1 prohibits fundamental separation.}

6. Therefore: \(\neg\Diamond C\) leads to a contradiction with A1.

7. Therefore: \(\Diamond C\).

\[
\boxed{\Diamond C}
\]

\textbf{This is A5 – but it is no longer an axiom, but a theorem.}

\subsection{A3 – No Absolute Nothing (Theorem)}\label{a3-no-absolute-nothing-theorem}

\textbf{Theorem:} From A1 follows \(\Box(N \rightarrow \neg\Diamond World)\).

\textbf{Proof:}

1. Absolute nothing \(N\) would be a \textbf{fundamentally separated realm} from reality.

2. \textbf{A1 prohibits fundamental separation.}

3. Therefore: \(N \rightarrow \neg\Diamond World\).

4. With Necessitation: \(\Box(N \rightarrow \neg\Diamond World)\).

\[
\boxed{\Box(N \rightarrow \neg\Diamond World)}
\]

\textbf{This is A3 – but it is no longer an axiom, but a theorem.}

\subsection{A2 – Modal Realism (Theorem)}\label{a2-modal-realism-theorem}

\textbf{Theorem:} From A1, A4 and A5 follows \(\forall p(\Diamond p \rightarrow \exists w\, Realized(w,p))\).

\textbf{Proof:}

1. Consciousness (C) is the capacity to experience \textbf{differences}. (From the definition of consciousness.)

2. If there were an \textbf{unrealized possibility} – a possibility that is not realized in any world – then this possibility would be \textbf{not experienceable by consciousness}.

3. An unexperienceable possibility would be \textbf{fundamentally separated} from the world of consciousness.

4. \textbf{A1 prohibits fundamental separation.}

5. Therefore: There can be no unrealized possibilities.

6. Therefore: \(\forall p(\Diamond p \rightarrow \exists w\, Realized(w,p))\).

\[
\boxed{\forall p(\Diamond p \rightarrow \exists w\, Realized(w,p))}
\]

\textbf{This is A2 – but it is no longer an axiom, but a theorem.}

\subsection{A6 – Origin as Ontological Embedding Condition (Definition + A10)}\label{a6-origin-as-ontological-embedding-condition-definition-a10}

\textbf{Theorem:} From the definition of U and the openness of T follows \(\Box(\neg Nec(p) \rightarrow \exists g\, Ground(g,p))\).

\textbf{Proof:}

1. Every non-necessary reality is \textbf{contingent} – it could also not exist.

2. If it exists, it needs a \textbf{ground} – otherwise it would be groundless.

3. The openness of T (as an open interval) guarantees that there is a \textbf{final ground}: the limit U.

4. Therefore: \(\Box(\neg Nec(p) \rightarrow \exists g\, Ground(g,p))\).

\[
\boxed{\Box(\neg Nec(p) \rightarrow \exists g\, Ground(g,p))}
\]

\textbf{This is A6 – but it is no longer an axiom, but follows from the definition of U and the openness of T.}

\subsection{A7 – Possibility of Complete Knowledge (Theorem)}\label{a7-possibility-of-complete-knowledge-theorem}

\textbf{Theorem:} From A1, A4 and A5 follows \(\Box(C \rightarrow \Diamond V_T)\).

\textbf{Proof:}

1. If consciousness exists (C), then it is \textbf{part of the one reality} (A1).

2. The one reality is \textbf{Totality T}.

3. If consciousness is part of T, then it \textbf{can} in principle know T – for there is no fundamental separation (A1) between knower and known.

4. Therefore: \(C \rightarrow \Diamond V_T\).

5. With Necessitation: \(\Box(C \rightarrow \Diamond V_T)\).

\[
\boxed{\Box(C \rightarrow \Diamond V_T)}
\]

\textbf{This is A7 – but it is no longer an axiom, but a theorem.}

\subsection{A8 – Reflexivity of Complete Knowledge (Theorem)}\label{a8-reflexivity-of-complete-knowledge-theorem}

\textbf{Theorem:} From A1 follows \(\Box(V_T(x) \rightarrow V_T(V_T(x)))\).

\textbf{Proof:}

1. If there is complete knowledge of T (\(V_T(x)\)), then this knowledge is \textbf{part of T} (for everything is part of T, A1).

2. If this knowledge is part of T, then it must also know \textbf{itself} – otherwise it would not be \textbf{complete}.

3. Therefore: \(V_T(x) \rightarrow V_T(V_T(x))\).

4. With Necessitation: \(\Box(V_T(x) \rightarrow V_T(V_T(x)))\).

\[
\boxed{\Box(V_T(x) \rightarrow V_T(V_T(x)))}
\]

\textbf{This is A8 – but it is no longer an axiom, but a theorem.}

\subsection{A9 – Identity of Knowledge and Object (Theorem)}\label{a9-identity-of-knowledge-and-object-theorem}

\textbf{Theorem:} From A1 follows \(\Box(V_T(x) \rightarrow T)\).

\textbf{Proof:}

1. Knowledge of T presupposes the existence of T. (This is tautological.)

2. Therefore: \(V_T(x) \rightarrow T\).

3. With Necessitation: \(\Box(V_T(x) \rightarrow T)\).

\[
\boxed{\Box(V_T(x) \rightarrow T)}
\]

\textbf{This is A9 – but it is no longer an axiom, but a theorem.}

\section{The Reductio ad absurdum}\label{the-reductio-ad-absurdum}

\subsection{Assumption}\label{assumption}

We assume that the Trinity does not exist:

\[
\neg Tr \equiv \neg(U \land T \land S)
\]

By de Morgan:

\[
\neg U \lor \neg T \lor \neg S
\]

We must show that each of the three cases leads to a contradiction.

\subsection{Case 1: Totality without Origin (\(T \land \neg U\))}\label{case-1-totality-without-origin-t-land-neg-u}

\textbf{Assumption:} \(T \land \neg U\)

\textbf{Proof of contradiction:}

1. T exists. So there is at least one state within the open interval.

2. Since T is an \textbf{open interval}, it has \textbf{no smallest element}. Every element has a "before."

3. But the openness of T \textbf{guarantees} the existence of the lower limit U (by definition).

4. If U does not exist (\(\neg U\)), then there is no lower limit.

5. Then T would no longer be an open interval – it would be either closed or infinite without a boundary.

6. \textbf{Contradiction to the definition of T.}

7. Therefore: \(T \land \neg U \rightarrow \bot\).

\[
\boxed{\Box(T \rightarrow U)}
\]

\subsection{Case 2: Origin without Self-Knowledge (\(U \land \neg S\))}\label{case-2-origin-without-self-knowledge-u-land-neg-s}

\textbf{Assumption:} \(U \land \neg S\)

\textbf{Proof of contradiction:}

1. U exists. This means: There is a lower limit of Totality.

2. From A5 (derived) it follows: \(\Diamond C\) – consciousness is possible.

3. From A2 (derived) it follows: \(\Diamond C \rightarrow \exists w\, Realized(w,C)\) – there is a world with consciousness.

4. From A7 (derived) it follows: \(C \rightarrow \Diamond V_T\) – complete knowledge is possible.

5. From A2 (derived) it follows: \(\Diamond V_T \rightarrow \exists w\, Realized(w,V_T)\) – there is a world with complete knowledge.

6. From A8 (derived) it follows: \(V_T(x) \rightarrow V_T(V_T(x))\) – complete knowledge knows itself. That is precisely S.

7. Therefore: \(U \rightarrow S\).

8. \textbf{Contradiction to the assumption} \(\neg S\).

\[
\boxed{\Box(U \rightarrow S)}
\]

\subsection{Case 3: Self-Knowledge without Totality (\(S \land \neg T\))}\label{case-3-self-knowledge-without-totality-s-land-neg-t}

\textbf{Assumption:} \(S \land \neg T\)

\textbf{Proof of contradiction:}

1. S exists. S is defined as "complete self-knowledge of Totality."

2. If S exists, then there is an \textbf{act of knowing} directed at T.

3. If T does not exist (\(\neg T\)), then S is \textbf{knowledge without an object}.

4. From A9 (derived) it follows: \(V_T(x) \rightarrow T\) – knowledge of T presupposes T.

5. Therefore: \(S \rightarrow T\).

6. \textbf{Contradiction to the assumption} \(\neg T\).

\[
\boxed{\Box(S \rightarrow T)}
\]

\section{The Synthetic Conclusion}\label{the-synthetic-conclusion}

From the three cases we have derived:

\[
\Box(T \rightarrow U), \quad \Box(U \rightarrow S), \quad \Box(S \rightarrow T)
\]

In S5, it follows:

\[
\Box(U \leftrightarrow T) \land \Box(T \leftrightarrow S) \land \Box(S \leftrightarrow U)
\]

With A4 (existence of a world) and the definition of T (as an open interval) it follows:

\[
\exists T \rightarrow \exists U \land \exists S
\]

Thus:

\[
\boxed{U \land T \land S}
\]

And with Necessitation:

\[
\boxed{\Box(U \land T \land S)}
\]

\textbf{The Trinity exists necessarily.}

\section{What Has Been Proved?}\label{what-has-been-proved}

\subsection{What the proof does not show}\label{what-the-proof-does-not-show}

\begin{itemize}
\item The existence of a personal God
\item A substantial threefoldness
\item A specific religious doctrine
\item A personal or emotional Trinity
\end{itemize}

\subsection{What the proof shows}\label{what-the-proof-shows}

\begin{itemize}
\item Under the axioms A1 and A4, the threefold limit structure of U, T, and S necessarily follows.
\item U, T, and S are not substances, but \textbf{limits} of an open interval.
\item The Trinity is the unity of these three limits – not three things, but three perspectives on the same reality.
\end{itemize}

\subsection{The three limits as perspectives}\label{the-three-limits-as-perspectives}

\begin{itemize}
\item \textbf{U (Origin)} – the perspective of "Whence?" – the lower limit that comes closest to nothing, but is itself not nothing.
\item \textbf{T (Totality)} – the perspective of "What is?" – the open interval of all realized states.
\item \textbf{S (Self-Knowledge)} – the perspective of "Who knows?" – the upper limit of complete self-transparency.
\end{itemize}

\section{Countermodels}\label{countermodels}

The Trinity can be avoided if at least one of the two axioms is abandoned:

\begin{tabular}{lll}
\textbf{Abandoned Axiom} & \textbf{Countermodel} & \textbf{Consequence} \\
A1 (Monism) & Dualism (Descartes) & S can exist without T; knowledge and object are separated \\
A4 (Existence of a World) & Nihilism & There is no world; the entire ontology is empty \\
\end{tabular}

\section{Comparison with Gödel's Ontological Argument}\label{comparison-with-goumldels-ontological-argument}

\subsection{Gödel's Argument (simplified)}\label{goumldels-argument-simplified}

\[
\text{Axioms about Positivity} \rightarrow \text{Necessary Existence of a Divine Being}
\]

\subsection{Comparison Table}\label{comparison-table}

\begin{tabular}{lll}
\textbf{Criterion} & \textbf{Gödel's Proof} & \textbf{This Proof} \\
Goal & Existence of a Being & Structural Necessity \\
Subject & Subject (God) & Limits (U, T, S) \\
Axioms & About "Positivity" & Monism + Existence of a World \\
Formal Rigor & High (but controversial) & High (explicitly verified) \\
Metaphysical Presuppositions & Strong (concept of positivity) & Minimal (only two axioms) \\
Philosophical Scope & Theological & Ontological-structural \\
Proximity to Classical Theology & Very high & Low (no person) \\
\end{tabular}

\section{Appendix: Complete Axioms and Theorems}\label{appendix-complete-axioms-and-theorems}

\subsection{Axioms (Complete)}\label{axioms-complete}

\[
\begin{aligned}
A1 &:= \Box\neg\exists x\exists y\, FundamentallySeparated(x,y) \\
A4 &:= \Diamond\exists w\, World(w)
\end{aligned}
\]

\subsection{Definitions}\label{definitions}

\[
\begin{aligned}
T &:= (U, S) \quad \text{(open interval)} \\
U &:= \lim_{x \to \inf} T \quad \text{(lower limit)} \\
S &:= \lim_{x \to \sup} T \quad \text{(upper limit)} \\
Tr &:= U \land T \land S
\end{aligned}
\]

\subsection{Derived Theorems}\label{derived-theorems}

\[
\begin{aligned}
& \Box(N \rightarrow \neg\Diamond World) \quad \text{(A3, from A1)} \\
& \Diamond C \quad \text{(A5, from A1 and A4)} \\
& \forall p(\Diamond p \rightarrow \exists w\, Realized(w,p)) \quad \text{(A2, from A1, A4, A5)} \\
& \Box(\neg Nec(p) \rightarrow \exists g\, Ground(g,p)) \quad \text{(A6, from definition of U)} \\
& \Box(C \rightarrow \Diamond V_T) \quad \text{(A7, from A1, A4, A5)} \\
& \Box(V_T(x) \rightarrow V_T(V_T(x))) \quad \text{(A8, from A1)} \\
& \Box(V_T(x) \rightarrow T) \quad \text{(A9, tautological)} \\
& \neg\exists infinite\, chain(g_1, g_2, ...) \quad \text{(A10, from openness of T)}
\end{aligned}
\]

\subsection{Derived Main Theorems}\label{derived-main-theorems}

\[
\begin{aligned}
& \Box(T \rightarrow U) \\
& \Box(U \rightarrow S) \\
& \Box(S \rightarrow T) \\
& \Rightarrow \Box(U \leftrightarrow T) \land \Box(T \leftrightarrow S) \land \Box(S \leftrightarrow U) \\
& \Rightarrow \Box(U \land T \land S) \\
& \Rightarrow \Box Tr
\end{aligned}
\]

\section{Conclusion}\label{conclusion}

The treatise has shown:

> Under the axioms of monism (A1) and the existence of a possible world (A4), the Trinity of Origin, Totality, and Self-Knowledge necessarily follows as a threefold limit structure of the one reality.

\textbf{The Trinity is not an additional entity, but a structural condition.}

It is what philosophy has sought since Plato and Plotinus, since Augustine and Hegel: the unity that carries its difference within itself, without falling into dualism or reductionist monism.

\begin{center}
\emph{This treatise was written in the spirit of strict modal logic, but in the language of philosophy – for truth requires both: the precision of the formula and the breadth of the concept.}
\end{center}

\end{document}